\chapter{Requirements, Ethics and Legal Position}
\label{chap:requirements}

This chapter states what the software had to do, the constraints it operated
under, and the ethical and legal position of the work.

\section{Functional requirements}
\label{sec:functional}

\begin{description}
  \item[F1 Reproduce the base method.] Run PerMFL unmodified on the authors'
    own dataset and compare against the published figure. Without this, no
    later result can be attributed to a change rather than to a broken
    implementation.
  \item[F2 Load intrusion detection data.] Parse CICIDS2017, NSL-KDD and
    TON-IoT into the four-tuple interface the existing loaders expose, so that
    every algorithm in the codebase runs on them without modification.
  \item[F3 Partition clients under controlled heterogeneity.] Support
    attack-domain, label-skew, Dirichlet and host-based partitioning, with the
    degree of skew settable and measurable.
  \item[F4 Report metrics appropriate to imbalanced detection.] Macro-averaged
    F1, precision, recall and false positive rate, per class and pooled, in
    addition to the accuracy the base implementation reports.
  \item[F5 Support repeated runs.] Vary model initialisation and partitioning
    by seed so that variation across runs can be measured.
  \item[F6 Derive teams from client updates.] Form teams during training from
    the client state rather than at load time, since a data-derived grouping
    cannot be computed before any client has trained.
  \item[F7 Evaluate on unseen classes.] Support evaluation of clients on
    classes absent from their training data.
\end{description}

\section{Non-functional requirements}
\label{sec:nonfunctional}

\begin{description}
  \item[N1 Reproducibility.] A run must be reproducible from a recorded seed
    and command. Package versions are pinned to the period of the base paper's
    publication.
  \item[N2 No information leakage.] No statistic fitted on evaluation data may
    influence training. Scaling parameters, clipping bounds and imputation
    values are fitted on the training split only.
  \item[N3 Commodity hardware.] Experiments must complete on a laptop CPU. No
    result may depend on access to a GPU cluster.
  \item[N4 Backward compatibility.] Every modification defaults to the shipped
    behaviour, so results obtained before a change remain reproducible after
    it.
  \item[N5 Traceability.] Each defect identified in the base implementation is
    recorded with a file and line reference and, where its effect is
    measurable, a measurement.
\end{description}

\section{Ethical position}
\label{sec:ethics}

\subsection{Data}

The project uses three public datasets released for research. CICIDS2017 and
NSL-KDD were produced by the Canadian Institute for Cybersecurity, TON-IoT by
UNSW Canberra. All three were captured in purpose-built testbeds using
synthetic user profiles and scripted attacks. No traffic from a real
organisation or an identifiable individual is involved, and no attempt is made
to re-identify any host.

The IP addresses in the captures belong to testbed machines. They appear in the
raw files and are used in one partitioning scheme to define client boundaries,
which reflects how a deployed system would be organised. They are removed
before the feature matrix is constructed, both to prevent the model learning
addresses instead of behaviour and because address-derived features do not
transfer outside the testbed.

\subsection{Ethical approval}

The project involves no human participants, no personal data as defined by UK
GDPR, and no interaction with live systems. Approval was obtained under
reference \missing{EthOS number}.

\subsection{Dual use}

Intrusion detection research is defensive, but any work that characterises
where a detector fails also indicates where it could be evaded. The findings
here concern the internal parameterisation of a federated learning method
rather than exploitable weaknesses in a deployed system, and the datasets are
public and well studied. The risk is low. No detection evasion technique is
developed or evaluated.

\subsection{Reporting}

Results are reported whether or not they support the hypothesis. The evaluation
in Chapter~\ref{chap:evaluation} includes a configuration in which the proposed
correction makes matters worse, and a claim from an earlier stage of the work
that later measurement contradicted. Both are retained.

\section{Legal position}
\label{sec:legal}

The base implementation is released by its authors under the MIT licence. This
work is a derivative and carries the same licence with attribution preserved.
Dataset licences permit academic use with citation, and all three are cited.
No dataset is redistributed; the repository contains loaders and a download
procedure.

\section{Scope}
\label{sec:scope}

Three exclusions are deliberate.

No comparison against SCMoE-PFL or CFMD-i is attempted. Neither publishes an
implementation, and comparing against a reimplementation would measure the
reimplementation. This also means no claim is made about relative computational
efficiency, since establishing one would require running both.

Differential privacy and secure aggregation are out of scope. The threat model
here is statistical heterogeneity, not an adversarial server.

Attack detection is treated as multi-class classification over labelled
traffic. Unsupervised anomaly detection, as in Fed-ANIDS, is a different
problem.
